\chapter{Conclusion and future work}
\label{ch:conclusion}

% Budget and sources: docs/thesis/outline.md. Every claim about behaviour points at a test, a figure or a number in data/.

\section{What the thesis shows}
\label{sec:conclusion-shows}

The thesis set out to show that payment disputes under different regulatory regimes do not need different
software: that one finite state machine, parameterised by the regime, can carry the lifecycle of a PSD2 card
dispute, a SEPA direct-debit claim, a Regulation~E debit-card error and a Regulation~Z billing error, with the
regime deciding which states exist, which deadlines run and which money moves. The engine built to test that claim
carries all four. Its fifteen states are shared, and each regime reaches the part of them its rules describe: six
under SEPA, eleven under PSD2, thirteen under Regulation~E and ten under Regulation~Z. None of the differences is
written as a test of a regime's name; each is a property of the rules, from the kind of money the customer receives
to whether a merchant may contest. The evaluation backs the claim with a test rather than an argument: every one of
the 64 transitions the four regimes accept is walked by a named scenario that ends in a closed dispute, and every
pair of state and event a regime forbids is refused (section~\ref{sec:eval-correctness}).

The rest of the work answers the question that follows from the first, which is whether a lifecycle held this way
can bear the weight of a real issuer's process. Deadlines are computed once, in the tenant's own business-day
calendar, and their status is derived on demand, so compliance is a distribution that can be read at any moment
rather than a report assembled afterwards. Money moves once: through a double-entry ledger whose suspense account
is zero whenever a dispute closes, and through a port to the bank's core whose duplicate detection makes a retry
after a timeout safe. Several institutions share one deployment, and the database rather than the code keeps
them apart; the actor on every event is the authenticated caller, and the analyst who investigated a dispute
cannot pay it out. These are the properties chapter~\ref{ch:requirements} put above any single feature, and each
is checked by a test against the real database rather than asserted
(sections~\ref{sec:eval-log} and~\ref{sec:eval-isolation}).

\section{What the decisions cost}
\label{sec:conclusion-costs}

Several decisions bought correctness at a price that the evaluation measured. Taking each authorization decision
inside the request's transaction, on facts already committed, doubled the median latency of a read from about two
to about four milliseconds; the cost is extra round trips to the database rather than the policy engine itself, and
the measurement shows where to recover it without moving the decision (section~\ref{sec:eval-latency}). Calling the
banking core inside the transaction makes the engine as slow as the slowest core it talks to, in exchange for never
recording a refund the core did not make. Keeping PostgreSQL as the only stateful dependency puts every
counter, session and outbox in the same database as the disputes; the measurements show its queries well under a
millisecond at this load, and the reversal is written down as a measurement to wait for rather than a preference.

Other decisions limited what the system claims rather than what it does. Fraud is scored by fixed rules whose
weights are judgement, not fitted to outcomes, and the score may only hold a credit until a person writes down why
it should go out; it never denies a dispute. The typed-decision model is offered to analysts only behind a
confidence gate, and its accuracy was measured on a synthetic set its author labelled. The banking core and the
card network are simulators behind ports. Each of these is stated in the text where it applies and collected in
section~\ref{sec:limitations}, so that what the thesis claims and what it measured do not drift apart.

Two parts were marked from the outset as the first to cut if time ran short: the fraud-scoring engine and the
analyst-facing workbench. Both were kept, in a form chosen to be defensible rather than impressive. The scoring is
readable and can only slow money down; the workbench never holds a token and shows only the actions the engine would
accept.

\section{Future work}
\label{sec:conclusion-future}

The most direct extensions deepen the lifecycle the thesis already models. The contested rounds of a chargeback
after the merchant's first answer, such as pre-arbitration and arbitration, are folded into one evidence step; they
would become states of their own, gated by the same properties, with deadlines from the card network's rules. The
statutory extensions of deadlines, such as the longer periods Regulation~E allows for new accounts or the holding
reply PSD2 permits in exceptional cases, would become clocks that an event may extend, with the extension recorded
beside the clock it moves. Both are additions to the rules rather than to the machine, which is the test the design
was meant to pass.

The next regulatory step is already visible. The European Union's Payment Services Regulation, provisionally agreed
at the end of 2025, extends the issuer's liability to some payments a customer was deceived into authorising, such as
those made to a fraudster impersonating the bank. Under it a dispute may arise from a payment that was authorised, the question becomes whether the bank's own fraud controls
met the expected standard, and the evidence comes from the bank rather than the merchant. That is a new regime in
the sense of this thesis: a row of rules, its clocks and its postings, and a test of whether the gates the four
existing regimes needed are enough for a fifth.

Integration would turn the simulators into adapters. A real core speaks the same instructions and receipts as the
simulated one, so the work is in the adapter and in observing a real core's behaviour when it is slow or silent,
which the port was designed for but has not met. The card network would receive the chargebacks that the engine now only
records, and its answers would become the events that the analyst now enters by hand.

Finally, the two measurements the running system already records would decide two of its open questions. The share
of the decision model's proposals that analysts keep, recorded whenever an analyst is shown one, is the honest
measure of whether the proposals help. And the outcomes of disputes, once there are enough of them, are
the labelled data a learned fraud model would need, against which the fixed rules could be compared rather than
replaced. Both are reasons the engine records what it records: the system was built so that its next decisions can
be taken on evidence, as its first ones were.
